\chapter{度量空间和赋范线性空间}
\section{度量空间的进一步例子}\label{sec:1.1}
\subsection{度量空间基础杂谈}
\prp{度量空间(距离空间)}{prp_2.1}{设 $X$ 是一个集合，若对于 $X$ 中任意两个元素 $x,y$，都有唯一确定的实数 $d(x,y)$ 与之对应，且这一对应关系满足下列条件：
\begin{enumerate}
    \item $d(x,y)\geq 0,d(x,y)=0$ 的充要条件为 $x=y$；
    \item $d(x,y)\leq d(x,z)+d(y,z)$，对任意 $z$ 都成立.
\end{enumerate}
则称 $d(x,y)$ 是 $x,y$ 之间的距离，称 $(X,d)$ 为度量空间或距离空间. $X$ 中的元素称为点，条件 2. 称为三点不等式.} 
 
$\left( \mathbb{R} ^n,d \right) $ 称为 n 维欧式空间，其中 $d$ 为欧几里得距离.

\ext{聚点、内点、界点}{ext_2.1}{
\begin{enumerate}
    \item {\color{magenta} 聚点}： $E\in \mathbb{R} ^n$，$P_0$ 为  $\mathbb{R} ^n$ 中一定点，如果 $P_0$ 的{\color{red} 任一邻域内都含有无穷多个属于 $E$ 的点}，称 $P_0$ 为 $E$ 的一个聚点.
    \item {\color{magenta} 内点}： 存在 $P_0$ 的某一邻域 $U(P_0)$ ，使 {\color{red} $U(P_0)\subset E$} ，则称 $P_0$ 为 $E$ 的内点.
    \item {\color{magenta} 界点}： $P_0$ 既非 $E$ 的内点又非 $E$ 的外点.
\end{enumerate}
}
\begin{prompt}
下面三个陈述等价：

\begin{enumerate}
    \item $P_0$ 为 $E$ 的一个聚点；
    \item  $P_0$ 的任一邻域内，至少含有一个属于 $E$ 而异于 $P_0$ 的点；
    \item 存在 $E$ 中互异的点所成列 $\left\{ P_n \right\} $ ，使 $P_n \rightarrow P_0 (n \rightarrow \infty)$.
\end{enumerate}
\end{prompt}

\ext{开核、导集、闭包}{ext_2.2}{
\begin{enumerate}
    \item {\color{magenta} 开核$\mathring{E}$}： $E$ 的{\color{red} 全体内点}组成的集合.
    \item {\color{magenta} 导集$E'$}： $E$ 的{\color{red} 全体聚点}组成的集合.
    \item {\color{magenta} 闭包$\overline{E}$}： {\color{red} $E\cup E'$} 称为 $E$ 的闭包， $\overline{E}=E\cup \partial E=\mathring{E}\cup \partial E=E'\cup \left\{ E\text{的孤立点} \right\} $
\end{enumerate}
}

\ext{开集、闭集、完备集}{ext_2.3}{
\begin{enumerate}
    \item {\color{magenta} 开集}：  $E\in \mathbb{R} ^n$，$E$ 中的每一点都是 $E$ 的{\color{red} 内点}，即 $E\subset \mathring{E},\text{或}E=\mathring{E}$.
    \item {\color{magenta} 闭集}：  $E\in \mathbb{R} ^n$，$E$ 中的每一个{\color{red} 聚点}都属于 $E$ ，即 $E'\subset \mathring{E},\text{或}\partial E\subset E$
    \item {\color{magenta} 自密集}：  $E\in \mathbb{R} ^n$，$E\subset E'$.
    \item {\color{magenta} 完备集}：  $E\in \mathbb{R} ^n$，$E=E'$.
\end{enumerate}}






\subsection{}














%--------------------------------%
%--------------------------------%
\section{度量空间中的极限，稠密集，可分空间}\label{sec:1.2}





%--------------------------------%
%--------------------------------%
\section{连续映射}\label{sec:1.3}




%--------------------------------%
%--------------------------------%
\section{柯西点列和完备度量空间}\label{sec:1.4}



%--------------------------------%
%--------------------------------%
\section{度量空间的完备化}\label{sec:1.5}



%--------------------------------%
%--------------------------------%
\section{压缩映射原理及其应用}\label{sec:1.6}




%--------------------------------%
%--------------------------------%
\section{线性空间}\label{sec:1.7}



%--------------------------------%
%--------------------------------%
\section{赋范线性空间和巴拿赫空间}\label{sec:1.8}